\section{Численный алгоритм}
\label{sec:algorithm}

\subsection{Прямая задача}
\label{sec:forward-method}
\draft{Вход: конфигурации, свойства среды, граничные условия и кинетические
константы. Результат: температурная история и зрелость витринита.
Последовательность решения:}
{\color{red!75!black}
\begin{enumerate}
 \item Построить расчётные сетки и найти начальное стационарное поле $T_0$.
       Начальные интегралы реакций равны нулю.
 \item Для каждой эпохи:
   \begin{enumerate}
    \item Задать движение сетки и моменты отложения материала.
    \item На каждом шаге до конца эпохи обновить геометрию, пористость
          и граничные данные; решить температурную задачу итерациями Пикара,
          выполняя линейное решение на каждой итерации.
    \item После сходимости температуры вычислить приращения интегралов и Ro;
          перейти к следующему временному шагу.
    \item При перестройке геометрии перенести состояние на следующую сетку.
   \end{enumerate}
 \item Завершить расчёт при $t=t_{\mathrm{final}}$.
\end{enumerate}}

\subsubsection{Устройство сетки}
\label{sec:mesh}
\draft{Область разбивается на колонки прямоугольных ячеек с независимыми
шагами $\Delta x$ и $\Delta z$. Сетка строится по контактам и движется
в соответствии с заданной историей среды. Несовпадающие уровни соседних
колонок образуют вертикальные подграни. Неизвестная $T_P$ относится к центру
ячейки $P$; рис.~\ref{fig:mesh} показывает сетку и участок у разлома.}
\reportfigure{02_mesh.png}
  {Модель 1: (а) --- сетка; (б), (в) --- участки со сдвигами ячеек.}
  {fig:mesh}

\subsubsection{Пространственный оператор и граничные потоки}
\draft{Применяется метод конечных объёмов~\cite{samarskii2003}.
Ячейки следуют движению материала, поэтому временное и конвективное слагаемые
учитываются материальной производной. Объёмные члены аппроксимируются
значениями в центре ячейки; $C_P=(\rho c)_P$:}
\begin{gather}
 \int_{P(t)}\rho c\left(\frac{\partial T}{\partial t}
              +\boldsymbol v\cdot\nabla T\right)\dd\Omega
      \approx C_P|P|\frac{\mathrm dT_P}{\mathrm dt},\qquad
 \int_{P(t)} A\dd\Omega\approx A_P|P|,\\
 \int_{P(t)}\nabla\cdot(\lambda\nabla T)\dd\Omega
       =-\sum_{\sigma\subset\partial P}F_{P\sigma},\qquad
 F_{P\sigma}=-\int_\sigma\lambda\nabla T\cdot\boldsymbol n_P\dd\sigma.
\end{gather}
\draft{Один внутренний поток входит в балансы соседних ячеек с противоположными
знаками: $F_{N\sigma}=-F_{P\sigma}$.
На горизонтальных гранях используется двухточечная аппроксимация
с учётом тепловых сопротивлений по обе стороны контакта;
$d_{P\sigma}$ --- расстояние от центра $P$ до грани $\sigma$:}
\begin{equation}
 F_{P\sigma}=\frac{|\sigma|(T_P-T_N)}
                    {d_{P\sigma}/\lambda_P+d_{N\sigma}/\lambda_N}.
\end{equation}
\draft{На вертикальной подграни глубины центров могут различаться;
разность $T_P-T_N$ тогда не определяет горизонтальную производную
(рис.~\ref{fig:flux}). Поток строится по нескольким соседним значениям T
с согласованием односторонних аппроксимаций.
Рабочая схема --- линейная MPFA; для сопоставления используется нелинейная NMPFA.
В линейной схеме коэффициенты потока зависят от геометрии и теплопроводности,
в нелинейной --- также от текущего приближения T:}
\begin{equation}
 F_{P\sigma}=\sum_{N\in\mathcal N_\sigma}a_{\sigma N}T_N+b_\sigma.
\end{equation}
\draft{Здесь $\mathcal N_\sigma$ --- соседние ячейки, участвующие в аппроксимации,
$b_\sigma$ учитывает граничные данные. Многоточечные схемы рассматриваются
в~\cite[гл.~2]{inmost2020}; точность выбранных схем проверяется
в разделе~\ref{sec:verification}.}
\reportfigure{03_flux.png}
  {Нормали к граням и направления между центрами соседних ячеек.}
  {fig:flux}

\draft{Граничные условия включаются в те же балансы потоков.
На горизонтальных участках поверхности задаётся температура, снизу ---
табличный поток на единицу горизонтальной проекции, на боковых границах ---
нулевой поток:}
\begin{equation}
 F_{P\sigma}^{\mathrm{surf}}=\frac{\lambda_P|\sigma|}{d_{P\sigma}}
                    (T_P-T_{\mathrm{surf}}),\qquad
 F_{P\sigma}^{\mathrm{basal}}=-q_{\mathrm{basal}}\Delta x_P,\qquad
 F_{P\sigma}^{\mathrm{side}}=0.
\end{equation}
\draft{На вертикальных ступеньках поверхности условие Дирихле учитывается
многоточечной аппроксимацией. Вертикальные ступеньки нижней границы
дополнительного притока не дают. В результате получается пространственный оператор:}
\begin{equation}
 C_P|P|\frac{\mathrm dT_P}{\mathrm dt}+\mathcal L_P(T)=0,\qquad
 \mathcal L_P(T)=\sum_{\sigma\subset\partial P}F_{P\sigma}(T)-A_P|P|.
\end{equation}

\subsubsection{Интегрирование по времени}
\draft{Используется неявный метод Эйлера. Шаги попадают в границы эпох;
геометрия и свойства берутся в новом моменте времени. Значения $T_P^n$
и $T_P^{n+1}$ относятся к одной движущейся ячейке:}
\begin{equation}
 C_P^{n+1}|P^{n+1}|\frac{T_P^{n+1}-T_P^n}{\Delta t}
       +\mathcal L_P^{n+1}(T^{n+1})=0.
 \label{eq:euler}
\end{equation}
\draft{Для начальной стационарной задачи временное слагаемое отсутствует:
$\mathcal L_P(T_0)=0$ при начальной геометрии и начальных условиях на границах.}

\subsubsection{Нелинейное решение}
\draft{Из-за зависимости свойств от T система~\eqref{eq:euler} нелинейна.
Реализован метод Пикара: на итерации $m$ фиксируются теплофизические
коэффициенты и, для NMPFA, веса потоков. Обозначив через $\boldsymbol R$
невязку системы~\eqref{eq:euler}, получим уравнение для поправки:}
\begin{equation}
 M(T^{(m)})\,\delta T=-\boldsymbol R(T^{(m)}),\qquad
 T^{(m+1)}=T^{(m)}+\delta T.
\end{equation}
\draft{Производные коэффициентов по T, необходимые для метода Ньютона,
в эту матрицу не включаются. Сходимость определяется абсолютной или
относительной нормой невязки температурной системы.}

\subsubsection{Линейное решение}
\draft{Используется GMRES(100)~\cite{saad1986} с предобуславливателем ILU(2)
и контролем предобусловленной невязки. Прямой решатель UMFPACK
служит вариантом для сопоставления. Выбор GMRES подтверждается
измерениями времени и различий T/Ro в разделе~\ref{sec:implementation}.}

\subsubsection{Вычисление зрелости после температурного шага}
\draft{После сходимости T вычисляются приращения интегралов реакций и Ro.
В каждой материальной точке на шаге температура считается линейной по времени;
интегрирование начинается с $t_* = \max(t^n,t_{\mathrm{dep}})$.
Если $t_*\geq t^{n+1}$, приращение равно нулю.
Обозначим $T_a=T(t_*)$, $T_b=T(t^{n+1})$, $\beta_i=E_i/R_g$.}
\begin{gather}
 G(\chi)=\frac{e^{-\chi}}{\chi}+\operatorname{Ei}(-\chi),\qquad
 \Phi_i(T)=\beta_iG(\beta_i/T),\\
 \Delta I_i=\frac{A_i(t^{n+1}-t_*)}{T_b-T_a}
             \bigl(\Phi_i(T_b)-\Phi_i(T_a)\bigr),\quad T_a\ne T_b,\\
 \Delta I_i=A_i(t^{n+1}-t_*)e^{-\beta_i/T_a},\quad T_a=T_b.
 \label{eq:ro-increment}
\end{gather}
\draft{Моменты отложения определены моделью~\eqref{eq:deposition}.
Размещение точек расчёта Ro и способ получения T в этих точках описаны
в разделе~\ref{sec:implementation}.}

\subsection{Обратная задача}
\label{sec:inverse-method}
\draft{Вход: наблюдения, оценки их погрешности и начальная история потока.
Результат: найденные значения параметров потока. Последовательность решения:}
{\color{red!75!black}
\begin{enumerate}
 \item Задать допустимое начальное приближение параметров.
 \item Для каждого пробного набора параметров выполнить полный прямой расчёт,
       включая стационарное начальное поле.
 \item Вычислить T и Ro в точках наблюдений, вектор невязок и функционал
       по формуле~\eqref{eq:objective}.
 \item Обновить параметры методом POUNDERS; повторять до выполнения
       критерия остановки.
\end{enumerate}}

\subsubsection{Параметризация потока}
\draft{Поток задан в 17 узлах на возрастах конфигураций и линейно
интерполируется между ними. Монотонность и границы обеспечиваются
параметрами $\eta_k\in[0,1]$:}
\begin{equation}
 q_0=q_{\min}+(q_{\max}-q_{\min})\eta_0,\qquad
 q_k=q_{\min}+(q_{k-1}-q_{\min})\eta_k,\quad k\geq1.
\end{equation}

\subsubsection{Безградиентная минимизация суммы квадратов}
\draft{Производные по параметрам не используются. POUNDERS
строит локальные модели отдельных компонент вектора невязок и на их основе
выбирает следующие параметры~\cite{wild2017,petscPounders}.
Каждое вычисление целевой функции требует одного прямого расчёта.}

\subsubsection{Остановка}
\draft{Критерии остановки: согласование всех наблюдений с погрешностью,
внутренний критерий оптимизатора, лимит итераций или число прямых вычислений.
Значения этих ограничений и фактический ход подбора приводятся
с результатами в разделе~\ref{sec:results}.}
